% §3 Exercises
%
% Cross-references and their resolution status:
%
% Already resolved (within chapter 3, drafted in Plan 3 Tasks 2-7):
%   \Cref{sec:and-under-noise}       -- §3.1 AND Under Noise
%   \Cref{eq:and-fnr}                -- §3.1 AND false-negative rate formula
%   \Cref{eq:and-fpr-general}        -- §3.1 AND false-positive rate formula
%   \Cref{sec:or-not}                -- §3.2 OR and NOT
%   \Cref{prop:boolean-closure}      -- §3.2 Boolean closure under noise
%   \Cref{sec:composition-theorem}   -- §3.3 Composition Theorem
%   \Cref{thm:composition}           -- §3.3 composition theorem statement
%   \Cref{eq:composition-theorem}    -- §3.3 eta_total formula
%   \Cref{sec:miller-rabin}          -- §3.4 Miller-Rabin walkthrough
%   \Cref{eq:mr-fpr}                 -- §3.4 single-round FPR bound
%   \Cref{eq:mr-k-round}             -- §3.4 k-round FPR bound
%   \Cref{sec:bridge-to-sets}        -- §3.5 Bridge to Sets
%
% Already resolved (chapter 4, drafted in Plan 1):
%   \Cref{sec:axiom-1}               -- chapter 4 §4.2
%   \Cref{ax:element-indep}          -- chapter 4 §4.2 Axiom 1
%   \Cref{sec:axiom-2}               -- chapter 4 §4.3
%   \Cref{ax:block-rates-indep}      -- chapter 4 §4.3 Axiom 2
%   \Cref{def:bernoulli-bool}        -- chapter 4 §4.5
%
% Unresolved forward refs (will resolve in Plan 4, Part I integration):
%   \Cref{ch:boolean-algebra}        -- chapter 3 itself (should resolve)
%   Chapter 16 reference is prose only (no \Cref) -- chapter 16 not yet drafted
%
% Difficulty-marker macros (not in alex.sty; defined locally if not yet defined):
\providecommand{\exstar}{$\star$}
\providecommand{\exstarstar}{$\star\star$}
\providecommand{\exstarstarstar}{$\star\star\star$}

\section*{Exercises}
\addcontentsline{toc}{section}{Exercises}

\begin{enumerate}

% -----------------------------------------------------------------------
% WARM-UP  (★ × 3)
% -----------------------------------------------------------------------

\item \exstar
\textbf{Deriving OR rates from AND via De Morgan.}
Let $b_1$ and $b_2$ be $\Bernoulli{\BoolType}$ instances with false-positive
rates $\fprate_1$, $\fprate_2$ and false-negative rates $\fnrate_1$, $\fnrate_2$.
Using De Morgan's law $b_1 \vee b_2 = \neg(\neg b_1 \wedge \neg b_2)$
and the NOT rule from \Cref{sec:or-not} (NOT swaps $\fprate$ and $\fnrate$),
derive closed-form expressions for the false-positive and false-negative
rates of $b_1 \vee b_2$.
Verify that your answer matches the direct truth-table derivation given
in \Cref{sec:or-not}.

\item \exstar
\textbf{FPR after a 3-stage AND chain (worked).}
Suppose three independent $\Bernoulli{\BoolType}$ instances $b_1$, $b_2$, $b_3$
each have false-positive rate $\fprate = 0.20$ and false-negative rate
$\fnrate = 0$.
\begin{enumerate}
  \item[(a)] Using the AND output-rate formula from \Cref{sec:and-under-noise},
    compute the false-positive rate $\fprate_{12}$ of $b_1 \wedge b_2$.
    (\textit{Hint:} with $\fnrate = 0$, several terms vanish; the
    only way to get a spurious $\top$ output is if both inputs report
    $\top$ when the true answer is $\bot$.)
  \item[(b)] Now compute the false-positive rate of
    $(b_1 \wedge b_2) \wedge b_3$, i.e., the full 3-stage AND.
  \item[(c)] Confirm that this agrees with the composition-theorem bound
    $\fprate_k \leq (\fprate)^k$ obtained by applying \Cref{thm:composition}
    in its special FNR-zero form, with $k = 3$ and the per-stage success
    probability $1 - \fprate = 0.80$ on negative inputs.
\end{enumerate}

\item \exstar
\textbf{Computing Miller-Rabin's $k$ for a target FPR.}
Miller-Rabin guarantees a per-round false-positive rate of at most
$\fprate_1 = 1/4$ (\Cref{eq:mr-fpr}).
\begin{enumerate}
  \item[(a)] Using the $k$-round bound $\fprate_k \leq (1/4)^k$
    (\Cref{eq:mr-k-round}), find the smallest integer $k$ such that
    $\fprate_k \leq 10^{-6}$.
    (\textit{Hint:} take logarithms base 10 on both sides.)
  \item[(b)] Find the smallest $k$ for $\fprate_k \leq 10^{-15}$
    (a bound suitable for cryptographic applications).
  \item[(c)] Suppose a weaker test had $\fprate_1 = 1/2$ per round.
    How many rounds would be needed to achieve $\fprate_k \leq 10^{-6}$?
    Comment on the practical difference.
\end{enumerate}

% -----------------------------------------------------------------------
% MEDIUM  (★★ × 5)
% -----------------------------------------------------------------------

\item \exstarstar
\textbf{Verifying the composition theorem on a 4-stage example.}
Consider a chain of four identity-with-error stages, each with per-stage
error rate $\eta_i = 0.10$.
\begin{enumerate}
  \item[(a)] Apply the composition theorem $\eta_{\mathrm{total}}
    = 1 - \prod_{i=1}^{k}(1 - \eta_i)$ (\Cref{eq:composition-theorem})
    with $k = 4$ and $\eta_i = 0.10$ to find $\eta_{\mathrm{total}}$.
  \item[(b)] Verify the same answer by building up stage by stage: compute
    $\eta_{1\oplus 2} = 1 - (1-\eta_1)(1-\eta_2)$, then
    $\eta_{1\oplus 2 \oplus 3}$, then the full 4-stage result.
    Confirm each step agrees with the cumulative formula.
  \item[(c)] Let $p_{\mathrm{correct}} = 1 - \eta_{\mathrm{total}}$ be the
    probability that all four stages produce a correct output simultaneously.
    Express $p_{\mathrm{correct}}$ as $(1 - \eta)^k$ and evaluate.
    Interpret the result: what fraction of queries emerge error-free from
    a four-stage pipeline with 10\% per-stage error?
\end{enumerate}

\item \exstarstar
\textbf{NOT is its own inverse on $\Bernoulli{\BoolType}$.}
Recall from \Cref{sec:or-not} that the NOT operation on a
$\Bernoulli{\BoolType}$ instance with rates $(\fprate, \fnrate)$ swaps
those rates, producing a new instance with rates $(\fnrate, \fprate)$.
\begin{enumerate}
  \item[(a)] Apply NOT twice to a $\Bernoulli{\BoolType}$ instance with
    rates $(\fprate, \fnrate)$. Show algebraically that the resulting
    rates are again $(\fprate, \fnrate)$.
  \item[(b)] Interpret the involution $(\mathrm{NOT})^2 = \mathrm{id}$ in
    terms of the channel matrix $Q$ from \Cref{def:bernoulli-bool}: what
    matrix operation corresponds to NOT, and why does applying it twice
    return the original matrix?
  \item[(c)] Does the same involution hold for AND and OR? That is, is
    there a notion of ``applying AND to itself'' that returns the identity?
    Explain briefly why the answer differs from the NOT case.
\end{enumerate}

\item \exstarstar
\textbf{Implement Miller-Rabin in 20 lines and check empirical FPR.}
Implement the Miller-Rabin primality test in the language of your choice.
The pseudocode in \Cref{sec:miller-rabin} fits in roughly 15 lines; a
full implementation with multi-round looping and result aggregation fits
in 20.
\begin{enumerate}
  \item[(a)] Test your implementation on the following inputs and report the
    outcome for each: $n = 2, 3, 4, 17, 100, 341, 561, 1009$.
    (Among these, 341 is a Fermat pseudoprime to base 2 but not a Miller-Rabin
    pseudoprime; 561 is a Carmichael number.
    Your implementation should correctly declare both composite with enough rounds.)
  \item[(b)] To estimate the empirical false-positive rate, generate a random
    sample of 1{,}000 composite odd integers in the range $[10^4, 10^6]$
    (e.g., products of two random primes). Run Miller-Rabin with $k = 1$
    round on each and record how many are incorrectly declared ``probably prime.''
    Compare your empirical rate to the theoretical bound of $1/4$.
  \item[(c)] Repeat part (b) with $k = 5$ and $k = 10$ rounds. Plot (or tabulate)
    empirical FPR vs.\ $k$. Does the empirical curve fall below the theoretical
    bound $(1/4)^k$ at each $k$?
\end{enumerate}

\item \exstarstar
\textbf{Prove the composition theorem assuming Axiom 1.}
The composition theorem (\Cref{thm:composition}) states that $k$ independent
identity-with-error stages in series produce an overall error rate
$\eta_{\mathrm{total}} = 1 - \prod_{i=1}^{k}(1 - \eta_i)$.
\begin{enumerate}
  \item[(a)] State precisely what independence assumption is needed at each
    stage of the chain for the per-stage probabilities to multiply. In the
    language of \Cref{sec:axiom-1}, what does it mean for chain-stage error
    events to be mutually independent?
  \item[(b)] Prove the formula by induction on $k$: the base case $k = 1$ is
    immediate; for the inductive step, use the chain rule and the
    independence of the $k$-th stage from all previous stages to factor
    the joint error probability.
  \item[(c)] As noted in \Cref{sec:composition-theorem}, this is
    \emph{chain-stage independence}, which is distinct from the element-wise
    Axiom~1 (\Cref{ax:element-indep}) even though both involve independence
    of error events. Write one paragraph explaining the operational difference:
    Axiom~1 governs \emph{which elements} are queried; chain-stage independence
    governs \emph{which stage} in the pipeline produces an error on a single
    query.
\end{enumerate}

\item \exstarstar
\textbf{Construct a 3-bit ANDed-confidence example.}
Let $b_1, b_2, b_3$ be three independent $\Bernoulli{\BoolType}$ instances
representing three noisy sensors, each observing whether a physical quantity
$x$ exceeds a threshold (true answer: $\top$). Suppose each sensor has
$\fprate = 0.15$ and $\fnrate = 0.05$.
\begin{enumerate}
  \item[(a)] Compute the false-positive and false-negative rates of the
    AND of all three sensors, $b_1 \wedge b_2 \wedge b_3$, by successive
    application of the AND formula from \Cref{sec:and-under-noise}.
    (\textit{Hint:} compute $b_1 \wedge b_2$ first, then AND with $b_3$.)
  \item[(b)] Compute the false-positive and false-negative rates of the
    OR of all three sensors, $b_1 \vee b_2 \vee b_3$, using the De Morgan
    derivation from \Cref{sec:or-not} or directly.
  \item[(c)] The AND system emphasizes low FPR; the OR system emphasizes
    low FNR. In numerical terms, which objective does the ANDed configuration
    better achieve compared with a single sensor? Which objective is
    \emph{worse} after ANDing? Justify your answer with the numbers from
    parts (a) and (b).
\end{enumerate}

% -----------------------------------------------------------------------
% CHALLENGING / EXTENSION  (★★★ × 3)
% -----------------------------------------------------------------------

\item \exstarstarstar
\textbf{Why the AND derivation in §3.1 tacitly assumes Axiom 2.}
Re-read the derivation of the AND output rates in \Cref{sec:and-under-noise}.
The derivation treats $\fprate_1$ and $\fprate_2$ as fixed constants when
writing the joint probability that both inputs produce a false $\top$.
\begin{enumerate}
  \item[(a)] Identify the exact expression in the derivation where the false-positive
    rates of the two inputs are multiplied together. Why does this step require
    that $\fprate_1$ and $\fprate_2$ be independent random quantities given the
    underlying set $A$?
    State which axiom (\Cref{ax:block-rates-indep}) licenses that factoring.
  \item[(b)] Suppose Axiom~2 (\Cref{ax:block-rates-indep}) fails for $b_1$ and $b_2$:
    specifically, suppose the false-positive rate of $b_2$ is a function of the
    false-positive rate of $b_1$, say $\fprate_2 = g(\fprate_1)$ for some
    nontrivial function $g$.
    Write a corrected expression for the false-positive rate of $b_1 \wedge b_2$
    in this dependent case (in terms of $\mathbb{E}[\fprate_1 \cdot g(\fprate_1)]$
    or a similar expectation).
  \item[(c)] Give a concrete scenario where this dependence arises naturally:
    for example, two sensors whose noise levels share a common environmental
    cause (correlated hardware failures, ambient temperature, etc.).
    Explain in two or three sentences why the simple product formula would
    overestimate the FPR reduction achieved by ANDing such sensors.
\end{enumerate}

\item \exstarstarstar
\textbf{De Morgan's laws hold for noisy Boolean gates.}
The standard De Morgan laws are $\neg(A \wedge B) = \neg A \vee \neg B$
and $\neg(A \vee B) = \neg A \wedge \neg B$.
Show that the same identities hold at the level of error rates for
$\Bernoulli{\BoolType}$ instances, using the rate-swapping NOT rule from
\Cref{sec:or-not} and the AND rate formulas from \Cref{sec:and-under-noise}.
\begin{enumerate}
  \item[(a)] Start from the OR-via-De-Morgan derivation in \Cref{sec:or-not}.
    Write out the step-by-step substitution: apply NOT to each input, apply
    AND, then apply NOT to the result. Track the false-positive and false-negative
    rates through each step.
  \item[(b)] Verify that the resulting rates for $\neg(b_1 \wedge b_2)$ equal
    those obtained by directly computing the rates of $(\neg b_1) \vee (\neg b_2)$,
    confirming the first De Morgan law.
  \item[(c)] In classical Boolean algebra, De Morgan's laws are
    \emph{definitional}: OR is \emph{defined} to be NOT-AND-NOT.
    Is the same true in the Bernoulli framework, or is the agreement of
    rates a theorem that requires proof?
    Identify any assumption (such as the closure property in
    \Cref{prop:boolean-closure}) that the argument needs.
\end{enumerate}

\item \exstarstarstar
\textbf{Constructing a system where the composition theorem fails.}
The composition theorem (\Cref{thm:composition}) requires that the per-stage
error events be mutually independent. This exercise shows that the theorem
can fail when that condition is violated -- a preview of the negative-space
discussion in Chapter~16.
\begin{enumerate}
  \item[(a)] Construct a two-stage identity-with-error system in which the
    error events at the two stages are perfectly positively correlated:
    whenever stage 1 errs, stage 2 errs with probability 1, and whenever
    stage 1 is correct, stage 2 is also correct.
    Let the per-stage error rate be $\eta = 0.10$.
    Compute the true end-to-end error rate $\eta_{\mathrm{total}}$.
  \item[(b)] Compare to the composition-theorem prediction
    $\eta_{\mathrm{total}}^{\mathrm{thm}} = 1 - (1-\eta)^2$.
    By how much does the actual error rate exceed the prediction, in absolute
    terms and as a ratio?
  \item[(c)] Generalise to $k$ perfectly correlated stages.
    Show that the actual error rate is $\eta$ (the per-stage rate) regardless
    of $k$, while the composition-theorem prediction approaches 0 as
    $k \to \infty$.
    In one paragraph, explain what this failure mode means for designing
    multi-stage pipelines in practice: when is the independence assumption
    most likely to be violated, and what structural property of the
    underlying system causes the correlation?
\end{enumerate}

\end{enumerate}
